\documentclass[12pt, a4paper]{article}
\usepackage{multicol}
\usepackage{geometry}
\usepackage{enumitem}
\usepackage{lipsum}
\usepackage{fontspec}

% Note: To use the BAHNSCHRIFT font from resources/fonts/, ensure the font is
% installed in the system or specify the correct path. For now, using a
% standard LaTeX font that will work out of the box.
\usepackage[T1]{fontenc}
\usepackage{lmodern}

% If you have the BAHNSCHRIFT.TTF font installed in your system, uncomment:
% \setmainfont{BAHNSCHRIFT}

\geometry{margin=1.5cm}

\setlength{\columnsep}{1.5cm}
\setlength{\columnseprule}{0.5pt}

\setlist[itemize]{nosep, topsep=0pt, leftmargin=1.5cm}

\begin{document}

\section*{\textbf{\large TRANSCONTINENTAL ACCENT GUIDE}}

\vspace{1cm}

\begin{multicols}{4}

\begin{minipage}{\columnwidth}
\textbf{1. Placement \& Posture}

\begin{itemize}
    \item Move the sound forward. Push resonance toward the front of your face and behind your teeth rather than deep in your throat.
    \item Hold a poised, brisk rhythm. Speak with an energetic, musical bounce—pitching slightly higher on key words and landing firmly at the end of thoughts.
\end{itemize}
\end{minipage}

\begin{minipage}{\columnwidth}
\textbf{2. Consonant Rules}

\begin{itemize}
    \item \textbf{Drop your trailing R's:} If an "R" comes after a vowel or at the end of a word, turn it into an "uh" or "ah" sound.
    \item \textbf{Examples:} Car $\rightarrow$ Cah, father $\rightarrow$ Fah-thuh
    \item \textbf{Tap your middle T's:} Never soften a "T" into a "D" sound.
    \item \textbf{Examples:} Butter $\rightarrow$ Buh-tuh, water $\rightarrow$ Wah-tuh
    \item \textbf{Add a "Y" glide before U's:} Insert subtle "y" sound after D, T, or N followed by long "U".
    \item \textbf{Examples:} Duty $\rightarrow$ Dyoo-tee, tune $\rightarrow$ Tyoon
    \item \textbf{Blow air into "WH" words:} Pronounce "H" before "W" in what, where, whine.
    \item \textbf{Finish every word:} Never drop "G" in -ing endings.
    \item \textbf{Example:} stepping $\rightarrow$ step-ping
\end{itemize}
\end{minipage}

\begin{minipage}{\columnwidth}
\textbf{3. Vowel Rules}

\begin{itemize}
    \item \textbf{Broaden your short A's:} In words like bath, can't, and dance, open jaw toward softer, rounder "ah".
    \item \textbf{Example:} can't $\rightarrow$ cahn't
    \item \textbf{Relax final E's:} Words ending in "Y" shouldn't end with tight "ee".
    \item \textbf{Example:} happy $\rightarrow$ hap-pih
    \item \textbf{Keep short O's open:} Unlike modern British English, keep words like not and hot open.
    \item \textbf{Examples:} not $\rightarrow$ naht, hot $\rightarrow$ haht
\end{itemize}

\vspace{0.5cm}

\textbf{4. Practice Drills}

\begin{itemize}
    \item \textit{"Could you pahss the buh-tuh, Wah-tuh?"}
    \item \textit{"It's Tom's dyoo-tee to do the lighting on Tyooz-day."}
    \item \textit{"I simply cahn't imagine wh-at the maht-tuh could be!"}
\end{itemize}

\end{minipage}

\begin{minipage}{\columnwidth}
\vspace{0.5cm}

\textbf{\large Keep Practicing}

If you'd like a customized breakdown:

\begin{itemize}
    \item Tell me: A specific line or monologue you're working on
    \item Tell me: Whether you're starting from a British or American baseline
\end{itemize}

I can write out a customized phonetic breakdown for you.

\begin{center}
\textit{...and remember, it's a theatrical hybrid designed for maximum clarity on old microphones.}
\end{center}
\end{minipage}

\end{multicols}

\end{document}